% GENERATED FILE. Edit canonical lessons, then run npm run generate:books.
% canonical-source-hash: fnv1a64:f61e908ae85513d5
% canonical-lessons: GU-C341-dalil, GU-C341-carca, GU-C341-muddo, GU-C341-sandarbh, GU-C341-uthavvu

\chapter{Making a Point}
\label{ch:gu-making-a-point}
\hlchaptermodality{341}

\begin{chapteropening}
\textbf{By the end of this chapter:} \emph{I can make a point in a discussion and back it up.}

Complete the last lesson of chapter 341: I can make a point in a discussion and back it up.
\end{chapteropening}

\section[\texorpdfstring{\gu{દલીલ} (dalīl)}{દલીલ (dalīl)}]{\gu{દલીલ} (dalīl) — a line of argument; a reason put forward}

\label{lesson:GU-C341-dalil}



\begin{quote}
Before the new one: say the Gujarati for an excuse, then the Gujarati for to accomplish.
\end{quote}

\subsection*{You'll want to know: \gu{દલીલ}}
\textbf{\gu{દલીલ}} — \emph{dalīl} — \textquotedblleft{}a line of argument; a reason put forward\textquotedblright{}.

\begin{quote}
\gu{મીરાની} \gu{દલીલ} \gu{મજબૂત} \gu{હતી}.
\end{quote}

\begin{quote}
\emph{Mīrānī dalīl majbūt hatī.}
\end{quote}

\begin{quote}
Mira's argument was a strong one.
\end{quote}

A feminine noun.

\begin{cousinweb}
\textbf{\gu{દલીલ}} (\emph{dalīl}) is usually traced to Arabic \emph{dalīl}, \textquotedblleft{}a guide; a proof\textquotedblright{}, through Persian. Gujarati uses it for a reason put forward in an argument.
\end{cousinweb}

\begin{grammarlens}[title={What you've built}]
One more word for giving reasons and opinions.
\end{grammarlens}

\subsection*{Your turn}
\begin{itemize}
\raggedright
  \item \emph{Say it:} \emph{dalīl}
  \item \emph{Say it:} \emph{dalīl}, once more
  \item \emph{Say it:} \emph{Mīrānī dalīl majbūt hatī.} — Mira's argument was a strong one.
  \item \emph{Recall:} say the Gujarati for an excuse, then the Gujarati for to accomplish, then say \emph{dalīl} again
\end{itemize}

\begin{tcolorbox}[breakable,colback=teal!4,colframe=teal!35!black,title={Before you move on}]
What does \gu{દલીલ} mean? (\textquotedblleft{}A line of argument; a reason put forward\textquotedblright{}.) Say it once more.
\end{tcolorbox}
\section[\texorpdfstring{\gu{ચર્ચા} (carcā)}{ચર્ચા (carcā)}]{\gu{ચર્ચા} (carcā) — a discussion}

\label{lesson:GU-C341-carca}



\begin{quote}
Before the new one: say the Gujarati for to accomplish, then the Gujarati for a line of argument.
\end{quote}

\subsection*{You'll want to know: \gu{ચર્ચા}}
\textbf{\gu{ચર્ચા}} — \emph{carcā} — \textquotedblleft{}a discussion\textquotedblright{}.

\begin{quote}
\gu{બેઠકમાં} \gu{નવા} \gu{નિયમ} \gu{પર} \gu{લાંબી} \gu{ચર્ચા} \gu{થઈ}.
\end{quote}

\begin{quote}
\emph{beṭhakmā̃ navā niyam par lā̃bī carcā thaī.}
\end{quote}

\begin{quote}
There was a long discussion about the new rule at the meeting.
\end{quote}

A feminine noun.

\begin{cousinweb}
\textbf{\gu{ચર્ચા}} (\emph{carcā}) is a tatsama, borrowed from Sanskrit \emph{carcā}, \textquotedblleft{}repetition, going over a thing again\textquotedblright{} (Monier-Williams); \textquotedblleft{}a discussion\textquotedblright{} is the sense Gujarati uses.
\end{cousinweb}

\begin{grammarlens}[title={What you've built}]
One more word for giving reasons and opinions.
\end{grammarlens}

\subsection*{Your turn}
\begin{itemize}
\raggedright
  \item \emph{Say it:} \emph{carcā}
  \item \emph{Say it:} \emph{carcā}, once more
  \item \emph{Say it:} \emph{beṭhakmā̃ navā niyam par lā̃bī carcā thaī.} — There was a long discussion about the new rule at the meeting.
  \item \emph{Recall:} say the Gujarati for to accomplish, then the Gujarati for a line of argument, then say \emph{carcā} again
\end{itemize}

\begin{tcolorbox}[breakable,colback=teal!4,colframe=teal!35!black,title={Before you move on}]
What does \gu{ચર્ચા} mean? (\textquotedblleft{}A discussion\textquotedblright{}.) Say it once more.
\end{tcolorbox}
\section[\texorpdfstring{\gu{મુદ્દો} (muddo)}{મુદ્દો (muddo)}]{\gu{મુદ્દો} (muddo) — a point at issue}

\label{lesson:GU-C341-muddo}



\begin{quote}
Before the new one: say the Gujarati for a line of argument, then the Gujarati for a discussion.
\end{quote}

\subsection*{You'll want to know: \gu{મુદ્દો}}
\textbf{\gu{મુદ્દો}} — \emph{muddo} — \textquotedblleft{}a point at issue\textquotedblright{}.

\begin{quote}
\gu{તમારો} \gu{મુદ્દો} \gu{સાચો} \gu{છે}.
\end{quote}

\begin{quote}
\emph{tamāro muddo sāco che.}
\end{quote}

\begin{quote}
Your point is right.
\end{quote}

A masculine noun.

\begin{cousinweb}
\textbf{\gu{મુદ્દો}} (\emph{muddo}) is usually traced to Arabic \emph{muddaʿā}, \textquotedblleft{}a thing claimed\textquotedblright{}, through Persian; in Gujarati it is the point a discussion turns on.
\end{cousinweb}

\begin{grammarlens}[title={What you've built}]
One more word for giving reasons and opinions.
\end{grammarlens}

\subsection*{Your turn}
\begin{itemize}
\raggedright
  \item \emph{Say it:} \emph{muddo}
  \item \emph{Say it:} \emph{muddo}, once more
  \item \emph{Say it:} \emph{tamāro muddo sāco che.} — Your point is right.
  \item \emph{Recall:} say the Gujarati for a line of argument, then the Gujarati for a discussion, then say \emph{muddo} again
\end{itemize}

\begin{tcolorbox}[breakable,colback=teal!4,colframe=teal!35!black,title={Before you move on}]
What does \gu{મુદ્દો} mean? (\textquotedblleft{}A point at issue\textquotedblright{}.) Say it once more.
\end{tcolorbox}
\section[\texorpdfstring{\gu{સંદર્ભ} (sandarbh)}{સંદર્ભ (sandarbh)}]{\gu{સંદર્ભ} (sandarbh) — a context}

\label{lesson:GU-C341-sandarbh}



\begin{quote}
Before the new one: say the Gujarati for a discussion, then the Gujarati for a point at issue.
\end{quote}

\subsection*{You'll want to know: \gu{સંદર્ભ}}
\textbf{\gu{સંદર્ભ}} — \emph{sandarbh} — \textquotedblleft{}a context\textquotedblright{}.

\begin{quote}
\gu{આ} \gu{શબ્દનો} \gu{અર્થ} \gu{સંદર્ભ} \gu{પર} \gu{આધાર} \gu{રાખે} \gu{છે}.
\end{quote}

\begin{quote}
\emph{ā śabdano arth sandarbh par ādhār rākhe che.}
\end{quote}

\begin{quote}
The meaning of this word depends on the context.
\end{quote}

A masculine noun.

\begin{cousinweb}
\textbf{\gu{સંદર્ભ}} (\emph{sandarbh}) is a tatsama, borrowed from Sanskrit \emph{saṃdarbha}, \textquotedblleft{}a stringing together\textquotedblright{} (Monier-Williams); Gujarati uses it for the setting that gives words their sense.
\end{cousinweb}

\begin{grammarlens}[title={What you've built}]
One more word for giving reasons and opinions.
\end{grammarlens}

\subsection*{Your turn}
\begin{itemize}
\raggedright
  \item \emph{Say it:} \emph{sandarbh}
  \item \emph{Say it:} \emph{sandarbh}, once more
  \item \emph{Say it:} \emph{ā śabdano arth sandarbh par ādhār rākhe che.} — The meaning of this word depends on the context.
  \item \emph{Recall:} say the Gujarati for a discussion, then the Gujarati for a point at issue, then say \emph{sandarbh} again
\end{itemize}

\begin{tcolorbox}[breakable,colback=teal!4,colframe=teal!35!black,title={Before you move on}]
What does \gu{સંદર્ભ} mean? (\textquotedblleft{}A context\textquotedblright{}.) Say it once more.
\end{tcolorbox}
\section[\texorpdfstring{\gu{ઉઠાવવું} (uṭhāvvũ)}{ઉઠાવવું (uṭhāvvũ)}]{\gu{ઉઠાવવું} (uṭhāvvũ) — to raise a point; to raise a question}

\label{lesson:GU-C341-uthavvu}



\begin{quote}
Before the new one: say the Gujarati for a point at issue, then the Gujarati for a context.
\end{quote}

\subsection*{You'll want to know: \gu{ઉઠાવવું}}
\textbf{\gu{ઉઠાવવું}} — \emph{uṭhāvvũ} — \textquotedblleft{}to raise a point; to raise a question\textquotedblright{}.

\begin{quote}
\gu{મીરાએ} \gu{બેઠકમાં} \gu{પગારનો} \gu{મુદ્દો} \gu{ઉઠાવ્યો}.
\end{quote}

\begin{quote}
\emph{Mīrāe beṭhakmā̃ pagārno muddo uṭhāvyo.}
\end{quote}

\begin{quote}
At the meeting Mira raised the question of pay.
\end{quote}

Past, with \gu{તેણે} (he or she …): \gu{તેણે} \gu{એક} \gu{પ્રશ્ન} \gu{ઉઠાવ્યો}, \emph{teṇe ek praśna uṭhāvyo}.

\begin{cousinweb}
\textbf{\gu{ઉઠાવવું}} (\emph{uṭhāvvũ}) is inherited from a form CDIAL reconstructs as \emph{ut-sthāpayati}, \textquotedblleft{}raises\textquotedblright{}; CDIAL lists the Gujarati verb (CDIAL 1903).
\end{cousinweb}

\begin{grammarlens}[title={What you've built}]
One more word for giving reasons and opinions.
\end{grammarlens}

\subsection*{Your turn}
\begin{itemize}
\raggedright
  \item \emph{Say it:} \emph{uṭhāvvũ}
  \item \emph{Say it:} \emph{uṭhāvvũ}, once more
  \item \emph{Say it:} \emph{Mīrāe beṭhakmā̃ pagārno muddo uṭhāvyo.} — At the meeting Mira raised the question of pay.
  \item \emph{Recall:} say the Gujarati for a point at issue, then the Gujarati for a context, then say \emph{uṭhāvvũ} again
\end{itemize}

\begin{tcolorbox}[breakable,colback=teal!4,colframe=teal!35!black,title={Before you move on}]
What does \gu{ઉઠાવવું} mean? (\textquotedblleft{}To raise a point; to raise a question\textquotedblright{}.) Say it once more.
\end{tcolorbox}
